\section{Pooled digits in the aligned self sized radix}
\label{sec:natural}
\subsection{Exact inherited interface and the changed rows}
All coordinates in this section range over $\N$. The input is the
Round~17 normalized raw quadratic graph for \file{strict_natural229}, with
334 variables and 284 rows. Its complete inherited ordinary proof is
\file{ONE_ROW_SEMANTIC_ANCHOR_PROOF.md}. The present transformations alter
only the eleven private modular remainder bounds and, in one variant,
the exact binomial output bound. The six exact power requests, semantic
Pell anchor at the derived label $J$, and simultaneous index certificate
are unchanged.

Use the inherited notation
\[
 t\ge12,\quad V_0=(F^2+1)^t,\quad K=x,\quad N=B^K,
 \quad q=(NS)^K.
\]
Here $S$ is one plus the sum of the nine positive independent CRT
representatives, not an individual stack representative. In genuine
completed witnesses,
\[
 F=t!,\qquad B=\binom{F^2+t}{t}.
\]
These two identities are not assumed prematurely in soundness. Write
\[
 D_0=q^t,\quad D_1=q^{t+1}=qD_0,\quad
 D_y=q^{y+1},\quad D_x=q^x.
\]
There are eleven modular remainders, organized as follows.
\begin{center}\small
\begin{tabular}{@{}lll@{}}\toprule
Remainders & Common denominator & Number\\\midrule
$r_0,r_C$ for the factorial digit and index column & $D_0$ & 2\\
$r_a,r_b,r_c,r_d,r_J$ & $D_1$ & 5\\
$r_{S_0},r_{S_1}$ for the two stack columns & $D_y$ & 2\\
$r_u,r_v$ for quotient columns & $D_x$ & 2\\\bottomrule
\end{tabular}
\end{center}
The old separate bounds are $D_i=r_i+s_i+1$ with $s_i\in\N$.
The exact factorial digit additionally satisfies $q=B+o+1$, with
$o\in\N$.

\subsection{A coefficient budget lemma}
\begin{lemma}\label{lem:digitbudget}
Let $q\ge2$, $k\ge1$, and
\[
 r_i=\sum_{j=0}^{k-1}c_{ij}q^j,\qquad c_{ij}\in\N.
\]
If $\sum_i c_{ij}\le q-1$ for every $j$, then
$\sum_i r_i\le q^k-1$. Hence
\begin{equation}\label{eq:poolgeneral}
 q^k=\sum_i r_i+s+1
\end{equation}
has a natural slack $s$. Conversely, for arbitrary natural $r_i,s$,
\eqref{eq:poolgeneral} implies $r_i<q^k$ for every $i$, and the old separate
slacks are restored by the nonnegative polynomial formulas
\[
 s_i=s+\sum_{h\ne i}r_h.
\]
\end{lemma}
\begin{proof}
Summing by positions gives
\[
 \sum_i r_i=\sum_{j=0}^{k-1}\Bigl(\sum_i c_{ij}\Bigr)q^j
 \le(q-1)\sum_{j=0}^{k-1}q^j=q^k-1.
\]
This proves existence of $s=q^k-1-\sum_i r_i\ge0$.
Conversely, each $r_i$ is a nonnegative summand of $q^k-s-1<q^k$.
The displayed restoration satisfies $r_i+s_i+1=q^k$ identically.
\end{proof}
The completeness condition bounds the sum of the coefficients, not merely
each remainder separately. Separate old inequalities alone would not
justify pooling.

\subsection{Uniform spare capacity in the actual radix}
Construct genuine witnesses in the inherited noncircular order. Begin with
a finite derivation and choose its padded horizon $t\ge12$. Fix
$F=t!$, $B$, $K=x$, and $N=B^K$; choose the nine positive independent CRT
representatives modulo $N^2$. A further lift of the dominant $a$
representative by $N^2$ preserves every required congruence and ensures
$S>N^2$. The inherited cofinal choice is retained, so the semantic anchor
still dominates the required indices. Every independent representative is
strictly less than $S$. The unchanged endpoint and ceiling arguments give
\[
 K=x\ge F+1\ge2,\qquad B>F\ge1,\qquad N\ge4.
\]
Consequently
\begin{equation}\label{eq:naturalradixmargin}
 q=N^KS^K\ge16S^K.
\end{equation}

For an ordinary representative $n<S$ and a requested lower digit position
$0\le j\le k\le K$,
\[
 \binom nj\le\max(1,n^K)<S^K.
\]
The same assertion for the derived label uses the inherited positive
arithmetic-graph majorant
\[
 J\le S^\ell,\qquad \ell=2^{22},\qquad
 \ell(t+1)<F<K.
\]
For $0\le j\le t+1$, this gives
$\binom Jj\le\max(1,J^j)<S^K$. For the factorial upper index,
the inherited elementary estimate yields
\[
 \binom{F^2+t}{j}\le V_0\quad(0\le j\le t),
 \qquad V_0<B^2\le N.
\]
Since $S>N^2$, both $V_0$ and $B$ are strictly below $S^K$.
Thus every low-position coefficient used in the eleven remainders, and
also the exact output $B$, is at most $S^K-1$.

These are completeness estimates, after genuine witnesses are fixed.
The unchanged common-root Pell congruence identifies $(q+1)^n$ modulo
the requisite digit modulus. Since the relevant binomial coefficients
are below $q$, the canonical remainders are exactly the low base-$q$
polynomials used above. This explanation introduces no exponentiation
operation into the emitted quadratic rows.

\subsection{Four independent pools}
Replace each group's old private bounds by one row
\begin{equation}\label{eq:naturalP4}
 D_G=\sum_{i\in G}r_i+s_G+1,\qquad s_G\in\N.
\end{equation}
The largest group has five remainders, so at every position the coefficient
sum is at most $5(S^K-1)<q$. Lemma~\ref{lem:digitbudget} supplies the four
new slacks for canonical witnesses.

For soundness, apply the converse part of the lemma independently to
each pool. Every deleted old slack is restored by
$s_i=s_G+\sum_{h\in G\setminus\{i\}}r_h\ge0$.
No binomial interpretation is needed in this direction. All other rows
and coordinates are unchanged, so a zero of the pooled raw graph yields a
zero of the exact inherited graph. For completeness, construct canonical
inherited witnesses first and then use the coefficient bounds above.
The raw graph loses seven variables and seven rows: its size becomes
327 variables and 277 rows.

\subsection{Weighted mixing of adjacent denominator levels}
The identity $D_1=qD_0$ lets the first two groups share a single bound.
Put $T=r_a+r_b+r_c+r_d+r_J$. Replace their seven old bounds by
\begin{equation}\label{eq:naturalP3}
 qD_0=q(r_0+r_C)+T+s+1,\qquad s\in\N.
\end{equation}
Keep the $D_y$ and $D_x$ pools from \eqref{eq:naturalP4}.
For soundness, $q\ge1$ and $D_0\ge1$ imply
\[
 qr_0<qD_0,\quad qr_C<qD_0,\quad
 r_a,r_b,r_c,r_d,r_J<D_1.
\]
Thus $r_0,r_C<D_0$, and all seven old slacks
$D_i-r_i-1$ are nonnegative integers. Their substitution restores the
old equations identically.

For completeness, multiplication by $q$ shifts each of the two smaller
remainders' coefficient arrays upward by one position. All terms lie below
position $t+1$. At most seven original coefficients contribute at any
position, so their sum is at most $7(S^K-1)<q$.
Lemma~\ref{lem:digitbudget} provides $s$. This raw graph has
326 variables and 276 rows.

The required positivity is available before decoding any binomial.
The retained frame and anchor first identify all six exact powers. The
unchanged ceiling equations imply $B\ge1$, $V_0\ge1$, $F\ge1$, $B>F$,
and $K\ge2$. Hence $N=B^K\ge4$ and $q=(NS)^K>0$.
This is precisely the inherited bootstrap; it uses neither a removed
remainder bound nor the separate exact digit output bound.

\subsection{Absorbing the exact binomial output bound}
A further compatible change replaces \eqref{eq:naturalP3} by
\begin{equation}\label{eq:naturalP3B}
 qD_0=q(r_0+r_C)+T+BD_0+s+1
\end{equation}
and deletes $q=B+o+1$ and its private $o$. In canonical witnesses,
$BD_0$ contributes the single digit $B$ at position $t$. At each position,
at most eight quantities bounded by $S^K-1$ contribute. Therefore their
sum is less than $8S^K\le q/2<q$, proving completeness by the same lemma.

For soundness, the nonnegative extra term $BD_0$ cannot weaken the
remainder inequalities already proved. The new row also gives
$BD_0<qD_0$. Since $D_0>0$, it follows that $B<q$.
Thus $o=q-B-1\in\N$ restores the deleted exact output bound.
The pre-decoding positivity argument in the preceding subsection remains
unchanged. This raw graph has 325 variables and 275 rows.
All displayed rows remain quadratic: $qr_i$ and $BD_0$ are products of
existing coordinates, not additional quantified gates.

\subsection{Exact projection to the emitted quartics}
Each selected endpoint is obtained from one of these raw graphs by a
finite sequence of monic polynomial eliminations. Every restored
coordinate is certified as a positive rational combination of squares of
affine rational forms plus a polynomial with nonnegative integer
coefficients. Such a formula is nonnegative on $\N^n$. Its exact equality
to the integer restoration polynomial is checked after expansion, so
integrality follows from the restoration polynomial itself.

The independent checker replays all substitutions and retained rows,
checks each nonnegativity certificate, verifies that residual degree never
exceeds two, and collects the entire final sum of squares. A changed
elimination order is supplied as an explicit permutation of the raw rows;
its permutation property and exact relation to the fixed pooled source
are checked. No heuristic search result is used as a substitute for the
restoration identities.

Composing the raw witness constructions with these certified projections
proves existence equivalence with the inherited natural graph and therefore
\eqref{eq:calibration}. Arbitrary old zeros need not extend to pooled zeros
with the same private remainders; the completeness argument constructs
canonical witnesses afresh from the encoded finite refutation. This is
existential equivalence, not an asserted bijection of every intermediate
zero set.

\subsection{Attribution and the precise new composition}
The incoming \file{conzfc_graph_digits_2026-10-24} package supplies the
four-pool digit-budget idea for an older semantic cap. Round~17 supplies
the aligned $K=x$ radix, factorial/product fusion, strict quotient ranges,
cofinal representative choice and literal parent graph. The new composition
uses the stronger margin \eqref{eq:naturalradixmargin}, weighted mixing of
adjacent denominator levels, and absorption of the exact output bound.
The retained contractor and independently implemented projection checker
are credited to the earlier package100 code. None of these attributions
alters the inherited logical calibration.
